% Generated draft: requires mathematical and source review.
\documentclass[11pt]{article}
\usepackage[margin=25mm]{geometry}
\usepackage{fontspec}
\setmainfont{Latin Modern Roman}
\usepackage{amsmath,amssymb,mathtools,unicode-math}
\setmathfont{Latin Modern Math}
\usepackage{enumitem,needspace,xurl,hyperref}
\setlength{\parindent}{0pt}
\setlength{\parskip}{5pt}
\setlength{\emergencystretch}{3em}
\newcommand{\N}{\mathbb{N}}
\newcommand{\Z}{\mathbb{Z}}
\newcommand{\Q}{\mathbb{Q}}
\newcommand{\R}{\mathbb{R}}
\newcommand{\C}{\mathbb{C}}
\newcommand{\F}{\mathbb{F}}
\newcommand{\A}{\mathbb{A}}
\newcommand{\PP}{\mathbb{P}}
\newcommand{\E}{\mathbb{E}}
\newcommand{\eps}{\varepsilon}
\newcommand{\dd}{\,\mathrm{d}}
\begin{document}
\section*{UnsolvedMath ID: 30001467}
\textbf{Source problem code:} catalogue entry
\subsection{Definitions and mathematical statement}
Let \(\mathbb{CP}^{2}\) be the complex projective plane,
\[
\mathbb{CP}^{2}=(\C^{3}\setminus\{0\})/\sim,
\qquad z\sim w\iff w=\lambda z\text{ for some }\lambda\in\C\setminus\{0\}.
\]
A group \(\Gamma\) is word-hyperbolic if it admits a finite generating set whose Cayley graph, with its graph metric, has \(\delta\)-thin geodesic triangles for some \(\delta\geq 0\). Its Gromov boundary \(\partial\Gamma\) is the space of equivalence classes of geodesic rays from a fixed vertex, where two rays are equivalent if their images have finite Hausdorff distance, equipped with the usual Gromov boundary topology.

Write \(X(\mathbb{CP}^{2})\) for the tree of manifolds named in the supplied source record. The record does not reproduce its defining construction or fully specify its orientation conventions; this notation therefore retains its source-dependent meaning. The supplied background explicitly distinguishes this space from an analogous tree involving both orientations of \(\mathbb{CP}^{2}\), so that analogous space must not be substituted for the target.

The question is whether there exist a finitely generated word-hyperbolic group \(\Gamma\) and a homeomorphism
\[
h\colon X(\mathbb{CP}^{2})\longrightarrow\partial\Gamma.
\]
A fully self-contained formulation still requires the precise definition of \(X(\mathbb{CP}^{2})\) from the cited source.
\subsection{Short English statement}
Does some word-hyperbolic group have Gromov boundary homeomorphic to the particular tree of complex projective planes denoted by \(X(\mathbb{CP}^{2})\) in the source? Its defining construction and orientation conventions require source verification. The supplied record labels this question open.
\subsection{Sources}
\begin{itemize}
\item Supplied local UnsolvedMath registry record, permanent ID \(30001467\), exact source problem code \texttt{catalogue entry}, in the Oberwolfach Reports Open Problems collection. The record associates the problem with 2010 and supplies no named proposer.
\item \emph{Non-positive Curvature and Geometric Structures in Group Theory} (2010), Oberwolfach Reports, pp.\ 1165--1224, DOI \texttt{10.4171/owr/2010/20}. Supplied source report and literature reference labeled ``original OWR report'': \url{https://doi.org/10.4171/owr/2010/20}.
\item Supplied literature reference labeled ``Swiatkowski, trees of manifolds'': \url{https://doi.org/10.2140/gt.2020.24.593}.
\item Literature assessment dated exactly \texttt{2026-08-21}: the supplied curation labels the problem open, reports that the primary boundary-realization paper explicitly leaves \(X(\mathbb{CP}^{2})\) as a question and discusses a cobordism obstruction, and reports realization of an analogous two-orientation tree rather than \(X(\mathbb{CP}^{2})\) itself. This is supplied dated curation, not an independently verified literature assessment; its status and cited claims require independent verification.
\end{itemize}
\end{document}
